\section{平展, 光滑与非分歧态射}\label{src:A.3}
称概形态射 $f:X\to Y$ 为
\begin{itemize}
\item \newterm{光滑}{smooth}, 如果 $f$ 局部有限表现且平坦, 并且对每个 $y\in Y$, 几何纤维 $X_{\overline{\kappa(y)}}=X\times_Y\Spec\overline{\kappa(y)}$ 正则;
\item \newterm{平展}{étale}, 如果 $f$ 光滑且相对维数为 $0$, 即对所有 $y\in Y$ 有 $\dim X_y=0$;
\item \newterm{非分歧}{unramified}, 如果 $f$ 局部有限型\footnote{我们遵循 \cite{GR71,SP} 的约定而非 \cite{EGA}, 因为仅要求 $f$ 局部有限型, 而不是局部有限表现.}, 且每条几何纤维都离散且约化, 即对所有 $y\in Y$, 有 $X_y\simeq\coprod_i\Spec K_i$, 其中每个 $K_i$ 都是 $\kappa(y)$ 的可分域扩张.
\end{itemize}
称概形态射 $f:X\to Y$ 在 $x\in X$ 处\textit{光滑} (相应地, \textit{平展}, \textit{非分歧}), 如果存在 $x$ 的开邻域 $U\subseteq X$, 使得 $f|_U:U\to Y$ 光滑 (相应地, 平展, 非分歧). 在 \S\bookref{A.3.5-subsection} 中, 我们将讨论局部完全交与 syntomic 态射. 有如下蕴含:
\[
\text{非分歧}\Longleftarrow\text{平展}\Longrightarrow\text{光滑}\Longrightarrow\text{syntomic}\Longrightarrow\text{fppf}\Longrightarrow\text{fpqc}.
\]

\subsection{等价条件}\label{src:A.3.1-subsection}
光滑, 平展与非分歧态射有许多等价刻画. 建立这些等价性需要相当多的工作; 易于理解的论述可参见 \cite[第 $1$ 章]{Mil80} 和 \cite[\S$4.3$]{Liu02}.
\begin{theorem}[光滑性的等价条件]\label{src:A.3.1}
设 $f:X\to Y$ 为概形 (相应地, 局部 Noether 概形) 之间的局部有限表现态射. 下列条件等价:
\begin{enumerate}[(1)]
\item $f$ 光滑;
\item (微分判别法) 对每个点 $x\in X$, $f$ 在 $x$ 处平坦, 且 $\mathcal O_{X,x}$-模 $\Omega_{X/S,x}$ 可以由至多 $\dim_x X_{f(x)}$ 个元素生成 (等价地, 它是秩为 $\dim_x X_{f(x)}$ 的自由模);
\item (无穷小提升判别法) 对每个核幂零的环满射 $A\to A_0$ (相应地, 满足 $\ker(A\to A_0)\simeq A/\mathfrak m_A$ 的局部 Artin 环满射) 和每个由实线箭头组成的交换图
\[
\begin{tikzcd}[column sep=3em,row sep=2em]
\Spec A_0\arrow[r]\arrow[d,hook]&X\arrow[d,"f"]\\
\Spec A\arrow[r]\arrow[ur,dashed]&Y
\end{tikzcd}
\]
存在补全该图的虚线箭头;
\item (Jacobi 判别法) 对每个点 $x\in X$, 存在 $f(x)$ 的仿射开邻域 $\Spec B$ 和 $x$ 的仿射开邻域 $\Spec A\subseteq f^{-1}(\Spec B)$, 以及 $A$-代数同构
\[
B\simeq A[x_1,\ldots,x_n]/(h_1,\ldots,h_r),
\]
其中 $h_1,\ldots,h_r\in A[x_1,\ldots,x_n]$, $r\leq n$, 并且对前 $r$ 个 $x_i$ 求偏导数得到的 Jacobi 矩阵的行列式 $\det(\frac{\delta h_j}{\delta x_i})_{1\leq i,j\leq r}\in B$ 为单位. (映射 $\Spec A\to\Spec B$ 称为\newterm{标准光滑态射}{standard smooth morphism}.)
\end{enumerate}
若 $X$ 和 $Y$ 还局部 Noether, 且 $x\in X$ 的像为 $y\in Y$, 满足 $\kappa(x)=\kappa(y)$, 则 $f:X\to Y$ 在 $x$ 处光滑, 当且仅当
\begin{enumerate}[(1)]\setcounter{enumi}{3}
\item 存在 $\widehat{\mathcal O}_{Y,y}$-代数同构 $\widehat{\mathcal O}_{X,x}\simeq\widehat{\mathcal O}_{Y,y}[x_1,\ldots,x_r]$.
\end{enumerate}
若 $X$ 和 $Y$ 还在代数闭域 $\Bbbk$ 上光滑, 则 $f$ 在 $x\in X(\Bbbk)$ 处光滑, 当且仅当
\begin{enumerate}[(1)]\setcounter{enumi}{4}
\item 切空间上的诱导映射 $T_{X,x}\to T_{Y,y}$ 满射.
\end{enumerate}
\end{theorem}
\begin{proof}
见 \cite[习题 II.$8.6$, 命题 III.$10.4$]{Har77}, \cite[$0.22.6.1$, IV$_4$.$17.5,14$]{EGA} 和 \cite[Tags 01V9, 02H6, 02HX]{SP}.
\end{proof}
称 $f:X\to Y$ 为\newterm{相对维数为 $n$ 的光滑态射}{smooth of relative dimension $n$}, 如果 $f$ 光滑, 且每条纤维都等维, 维数为 $n$; 等价地, $f$ 为 fppf 态射, 所有纤维都等维, 维数为 $n$, 且 $\Omega_{X/S}$ 为秩 $n$ 的局部自由层. 若仅有 $f$ 为 fppf 态射, 且 $\Omega_{X/S}$ 为维数 $d$ 的局部自由层, 则 $f$ 未必是相对维数为 $d$ 的光滑态射, 例如 $\Spec\mathbb F_p[x]\to\Spec\mathbb F_p[x^p]$.
\begin{theorem}[平展性的等价条件]\label{src:A.3.2}
设 $f:X\to Y$ 为概形 (相应地, 局部 Noether 概形) 之间的局部有限表现态射. 下列条件等价:
\begin{enumerate}[(1)]
\item $f$ 平展;
\item $f$ 光滑且 $\Omega_{X/Y}=0$;
\item $f$ 光滑且非分歧;
\item $f$ 平坦且非分歧;
\item (无穷小提升判别法) 对每个核幂零的环满射 $A\to A_0$ (相应地, 满足 $\ker(A\to A_0)\simeq A/\mathfrak m_A$ 的局部 Artin 环满射) 和每个由实线箭头组成的交换图
\[
\begin{tikzcd}[column sep=3em,row sep=2em]
\Spec A_0\arrow[r]\arrow[d,hook]&X\arrow[d,"f"]\\
\Spec A\arrow[r]\arrow[ur,dashed]&Y
\end{tikzcd}
\]
存在唯一的虚线箭头补全该图;
\item (Jacobi 判别法) 对每个点 $x\in X$, 存在 $f(x)$ 的仿射开邻域 $\Spec B$ 和 $x$ 的仿射开邻域 $\Spec A\subseteq f^{-1}(\Spec B)$, 以及 $A$-代数同构
\[
B\simeq A[x_1,\ldots,x_n]/(h_1,\ldots,h_n),
\]
其中 $h_1,\ldots,h_n\in A[x_1,\ldots,x_n]$, 并且行列式 $\det(\frac{\delta h_j}{\delta x_i})_{1\leq i,j\leq n}\in B$ 为单位. (映射 $\Spec A\to\Spec B$ 称为\newterm{标准平展态射}{standard étale morphism}.)
\end{enumerate}
若 $X$ 和 $Y$ 还局部 Noether, 且 $x\in X$ 的像为 $y\in Y$, 满足 $\kappa(x)=\kappa(y)$, 则 $f:X\to Y$ 在 $x$ 处光滑, 当且仅当
\begin{enumerate}[(1)]\setcounter{enumi}{5}
\item $\widehat{\mathcal O}_{X,x}\simeq\widehat{\mathcal O}_{Y,y}$.
\end{enumerate}
若 $X$ 和 $Y$ 还在代数闭域 $\Bbbk$ 上光滑, 则 $f$ 在 $x\in X(\Bbbk)$ 处平展, 当且仅当
\begin{enumerate}[(1)]\setcounter{enumi}{6}
\item 切空间上的诱导映射 $T_{X,x}\to T_{Y,y}$ 为同构.
\end{enumerate}
\end{theorem}
\begin{proof}
见 \cite[习题 III.$10.3$]{Har77}, \cite[IV$_4$.$17.14.1$--$2$, IV$_4$.$17.6.3$]{EGA} 和 \cite[Tags 02GH, 02HF]{SP}.
\end{proof}
\begin{theorem}[非分歧性的等价条件]\label{src:A.3.3}
设 $f:X\to Y$ 为局部有限型概形态射. 下列条件等价:
\begin{enumerate}[(1)]
\item $f$ 非分歧;
\item $\Omega_{X/Y}=0$;
\item 对角态射 $\Delta_f:X\to X\times_Y X$ 为开浸入;
\item (非分歧性的无穷小提升判别法) 对每个核幂零的环满射 $A\to A_0$ (相应地, 满足 $\ker(A\to A_0)\simeq A/\mathfrak m_A$ 的局部 Artin 环满射) 和每个由实线箭头组成的交换图
\[
\begin{tikzcd}[column sep=3em,row sep=2em]
\Spec A_0\arrow[r]\arrow[d,hook]&X\arrow[d,"f"]\\
\Spec A\arrow[r]\arrow[ur,dashed]&Y
\end{tikzcd}
\]
至多存在一条虚线箭头补全该图.
\end{enumerate}
若 $X$ 和 $Y$ 还局部 Noether, 且 $x\in X$ 的像为 $y\in Y$, 满足 $\kappa(x)=\kappa(y)$, 则 $f:X\to Y$ 在 $x$ 处光滑, 当且仅当
\begin{enumerate}[(1)]\setcounter{enumi}{3}
\item $\widehat{\mathcal O}_{Y,y}\to\widehat{\mathcal O}_{X,x}$ 满射.
\end{enumerate}
\end{theorem}
\begin{proof}
见 \cite[IV$_4$.$17.14.1$--$2$, IV$_4$.$17.6.3$]{EGA} 和 \cite[Tags 02G3, 02H7, 02GE]{SP}.
\end{proof}

\subsection{光滑, 平展与非分歧态射的平展局部结构}\label{src:A.3.2-subsection}
每个光滑态射在平展局部都是相对仿射空间.
\begin{proposition}[光滑态射的局部结构]\label{src:A.3.4}
概形态射 $X\to Y$ 在 $x\in X$ 处光滑, 当且仅当存在仿射开子概形 $\Spec A\subseteq X$ 和 $\Spec B\subseteq Y$, 其中 $x\in\Spec A$, 以及交换图
\[
\begin{tikzcd}[column sep=2.5em,row sep=2em]
X\arrow[d]&\Spec A\arrow[l,hook',"\mathrm{op}"']\arrow[d]\arrow[r,"\text{ét}"]&\mathbb A_B^n\arrow[dl]\\
Y&\Spec B\arrow[l,hook',"\mathrm{op}"']&
\end{tikzcd}
\]
其中 $\Spec A\to\mathbb A_B^n$ 平展.
\end{proposition}
\begin{proof}
见 \cite[Tag 039P]{SP} 和 \cite[IV$_4$.$17.11.4$]{EGA}.
\end{proof}
一个重要推论是, 光滑态射在平展局部有截面.
\begin{corollary}\label{src:A.3.5}
设 $f:X\to Y$ 为在 $x\in X$ 处光滑的概形态射. 则存在 $f(x)$ 的平展邻域 $Y'\to Y$, 使得 $X\times_Y Y'\to Y'$ 有截面.
\end{corollary}
\begin{proof}
应用上述命题后, 注意态射 $\mathbb A_B^n\to\Spec B$ 有零截面 $\Spec B\to\mathbb A_B^n$. 概形 $Y':=\Spec B\times_{\mathbb A_B^n}\Spec A$ 在 $Y$ 上平展, 且复合 $Y'\to\Spec A\hookrightarrow X$ 定义 $X\times_Y Y'\to Y'$ 的截面 $Y'\to X\times_Y Y'$.
\end{proof}
每个平展 (相应地, 非分歧) 态射在平展局部都是同构 (相应地, 闭浸入).
\begin{proposition}\label{src:A.3.6}
设 $f:X\to S$ 为在 $x\in X$ 处平展 (相应地, 非分歧) 的分离概形态射. 则存在平展邻域 $(U,u)\to(S,f(x))$ 以及有限不交并分解
\[
X_U=W\coprod\coprod_i V_i,
\]
使得每个 $V_i\to U$ 都是同构 (相应地, 闭浸入), 且纤维 $W_u$ 不含位于 $x$ 上方的点.
\end{proposition}
\begin{proof}
见 \cite[Tags 04HM, 04HG]{SP}.
\end{proof}
有时知道仿射概形之间的平展与光滑态射可以沿闭浸入提升, 会很方便. 对 syntomic 态射也成立 (见定义 \bookref{A.3.17}).
\begin{proposition}\label{src:A.3.7}
考虑由实线箭头组成的图
\[
\begin{tikzcd}[column sep=3em,row sep=2em]
\Spec A_0\arrow[d]\arrow[r,hook,dashed]&\Spec A\arrow[d,dashed]\\
\Spec B_0\arrow[r,hook]\arrow[ur,phantom,"\square" description]&\Spec B
\end{tikzcd}
\]
其中 $\Spec B\hookrightarrow\Spec B_0$ 为闭浸入. 若 $\Spec A_0\to\Spec B_0$ 平展 (相应地, 光滑, syntomic), 则存在平展 (相应地, 光滑, syntomic) 态射 $\Spec A\to\Spec B$, 使上述图为 Cartesian 图.
\end{proposition}
\begin{proof}
见 \cite[Tags 04D1, 07M8]{SP}.
\end{proof}

\subsection{进一步性质}\label{src:A.3.3-subsection}
\begin{proposition}[纤维平展性/光滑性/非分歧性判别法]\label{src:A.3.8}
考虑概形的图
\[
\begin{tikzcd}[column sep=2em,row sep=1.5em]
X\arrow[rr]\arrow[dr]&&Y\arrow[dl]\\&S&
\end{tikzcd}
\]
其中 $X\to S$ 和 $Y\to S$ 局部有限表现. 设 $x\in X$ 的像为 $s\in S$. 则
\begin{enumerate}[(1)]
\item $X\to Y$ 在 $x$ 处非分歧, 当且仅当 $X_s\to Y_s$ 在 $x$ 处非分歧;
\item 若 $X\to S$ 在 $x$ 处平坦, 则 $X\to Y$ 在 $x$ 处平展 (相应地, 光滑), 当且仅当 $X_s\to Y_s$ 在 $x$ 处平展 (相应地, 光滑).
\end{enumerate}
\end{proposition}
\begin{proof}
设 $y\in Y$ 为 $x\in X$ 的像. 第 (1) 条成立, 因为非分歧性定义为纤维上的条件, 而纤维 $X_y$ 与 $X_s\to Y_s$ 在 $y\in Y_s$ 上的纤维等同. 对于第 (2) 条的非平凡方向 ($\Leftarrow$), 纤维平坦性判别法 (\bookref{A.2.10}) 意味着 $X\to Y$ 在 $x$ 处平坦. 因而 $X\to Y$ 在 $x$ 处光滑 (相应地, 平展), 等价于 $X_y\to\Spec\kappa(y)$ 在 $x$ 处光滑 (相应地, 平展). 后一条件由 $X_s\to Y_s$ 在 $x$ 处光滑 (相应地, 平展) 成立.
\end{proof}
光滑性显然是源上的开条件; 当态射固有时, 它也是靶上的开条件.
\begin{corollary}\label{src:A.3.9}
若 $f:X\to Y$ 为固有, 平坦且局部有限表现态射, 则满足 $X_y\to\Spec\kappa(y)$ 光滑的点 $y\in Y$ 的集合为开集.
\end{corollary}
\begin{proof}
若 $y\in Y$ 满足 $X_y\to\Spec\kappa(y)$ 光滑, 则 $f:X\to Y$ 在 $X_y$ 的某个开邻域上光滑. 若 $Z\subset X$ 为 $f:X\to Y$ 不光滑的闭子集, 则 $f(Z)\subseteq Y$ 恰好是 $f$ 的纤维不光滑的点集. 由于 $f$ 固有, $f(Z)$ 闭.
\end{proof}
\begin{proposition}\label{src:A.3.10}
设 $X\to Y$ 为 Noether 概形之间的光滑态射. 对每个像为 $y\in Y$ 的点 $x\in X$, 有
\[
\dim_x(X)=\dim_y(Y)+\dim_x(X_y).
\]
\end{proposition}
\begin{proof}
见 \cite[Tag 0AFF]{SP}.
\end{proof}
\begin{proposition}\label{src:A.3.11}
若 $X\to Y$ 为有限平展态射, 则存在有限平展覆盖 $Y'\to Y$, 使得 $X\times_Y Y'\to Y'$ 为平凡覆盖, 即 $X\times_Y Y'$ 在 $Y'$ 上同构于 $\coprod_i Y'$.
\end{proposition}
\begin{proof}
可以假设 $X\to Y$ 的次数 $d$ 为常数. 概形
\[
(X/Y)^d=\underbrace{X\times_Y\cdots\times_Y X}_{d}
\]
表出函子 $\Sch/Y\to\Sets$, 它把 $Y$-概形 $T$ 映到 $X\times_Y T\to T$ 的 $d$ 个截面组成的集合. 每个两两对角态射 $(X/Y)^{d-1}\to(X/Y)^d$ 都是既开又闭的浸入; 令 $(X/Y)^d_0\subseteq(X/Y)^d$ 为所有两两对角线的补集. 投影态射 $(X/Y)^d_0\to Y$ 有限平展, 而函子描述给出 $X\times_Y(X/Y)^d_0\to(X/Y)^d_0$ 的 $d$ 个互不相交的截面.
\end{proof}
\begin{proposition}\label{src:A.3.12}
若概形的支配非分歧态射 $X\to Y$ 满足 $Y$ 正规且 $X$ 连通, 则它平展.
\end{proposition}
\begin{proof}
见 \cite[推论 I.$9.11$]{SGA1}.
\end{proof}
下面的结果常称为 ``Nagata--Zariski 纯性''.
\begin{proposition}[分支点集的纯性]\label{src:A.3.13}
设 $f:X\to Y$ 为整 Noether 概形之间的拟有限态射, $X$ 正规且 $Y$ 正则. 则 $X$ 中 $f$ 不平展的点集为空或余维为 $1$.
\end{proposition}
\begin{proof}
见 \cite{Zar58}, \cite{Nag59}, \cite[定理 X.$3.1$]{SGA1} 和 \cite[Tag 0BMB]{SP}.
\end{proof}

\subsection{Fitting 理想与奇异点集}\label{src:A.3.4-subsection}
Fitting 理想能够给出概形的奇异点集的概形论描述. 我们在结点的刻画 (\bookref{6.2.4}) 中使用 Fitting 理想. 关于 Fitting 理想的背景, 推荐 \cite[Tag 07Z6]{SP} 和 \cite[\S$20$]{Eis95}.

若 $R$ 为环, $M$ 为有限生成 $R$-模, 则 $M$ 的\newterm{第 $k$ 个 Fitting 理想}{$k$th fitting ideal} $\operatorname{Fit}_k(M)$ 是由定义呈示
\[
\bigoplus_{i\in I}R\xrightarrow{A}R^n\longrightarrow M\longrightarrow0
\]
的矩阵 $A$ 的 $(n-k)\times(n-k)$ 子式生成的理想. 当然, 当 $M$ 有限表现 (例如 $R$ 为 Noether 环) 时, 可假设左端项为有限自由模 $R^m$; 此时 $A$ 为 $m\times n$ 矩阵, $\operatorname{Fit}_k(M)$ 为有限生成理想. Fitting 理想与呈示的选取无关, 并定义递增理想序列
\[
0=\operatorname{Fit}_{-1}(M)\subseteq\operatorname{Fit}_0(M)\subseteq\operatorname{Fit}_1(M)\subseteq\cdots\subseteq R,
\]
使得若 $M$ 可以由 $k$ 个元素生成, 则 $\operatorname{Fit}_k(M)=R$. $R$-模 $M$ 为秩 $r$ 的局部自由模, 当且仅当 $\operatorname{Fit}_{r-1}(M)=0$ 和 $\operatorname{Fit}_r(M)=R$; 此时对所有 $k<r$ 有 $\operatorname{Fit}_k(M)=0$. 对环映射 $R\to S$, 有等同 $\operatorname{Fit}_k(M\otimes_R S)=\operatorname{Fit}_k(M)S$. 特别地, 对 $f\in R$, 有 $\operatorname{Fit}_k(M_f)=\operatorname{Fit}_k(M)_f$; 对素理想 $\mathfrak p\subseteq R$, 有 $\operatorname{Fit}_k(M_{\mathfrak p})=\operatorname{Fit}_k(M)_{\mathfrak p}$; 若 $R$ 为 Noether 局部环, 则 $\operatorname{Fit}_k(M)\otimes_R\widehat R=\operatorname{Fit}_k(\widehat M)$.

若 $X$ 为概形, $F$ 为 $X$ 上的有限型拟凝聚层, 则 $F$ 的\newterm{第 $k$ 个 Fitting 理想层}{$k$th fitting ideal sheaf} 为拟凝聚理想层 $\operatorname{Fit}_k(F)\subseteq\mathcal O_X$, 它由对仿射开子集 $U\subseteq X$ 的等式 $\Gamma(U,\operatorname{Fit}_k(F))=\operatorname{Fit}_k(\Gamma(F,U))$ 定义. Fitting 理想层给奇异点集赋予概形结构.
\begin{definition}\label{src:A.3.14}
若 $X$ 为域 $\Bbbk$ 上纯 $d$ 维的 Noether 概形, 则定义 $X$ 的\newterm{奇异点集}{singular locus} 为子概形 $\operatorname{Sing}(X):=V(\operatorname{Fit}_d(\Omega_{X/\Bbbk}))$, 它由 $\Omega_{X/\Bbbk}$ 的第 $d$ 个 Fitting 理想定义. 更一般地, 若 $X\to S$ 为 fppf 态射, 每条纤维都纯 $d$ 维, 则定义\newterm{相对奇异点集}{relative singular locus} 为子概形 $\operatorname{Sing}(X/S):=V(\operatorname{Fit}_d(\Omega_{X/S}))$.
\end{definition}
例如, 若 $X=\Spec\Bbbk[x_1,\ldots,x_n]/I$, $I=(f_1,\ldots,f_m)$, 则正合列 $I/I^2\to\Omega_{\mathbb A^n/\Bbbk}|_X\to\Omega_{X/\Bbbk}\to0$ 诱导消解
\[
\mathcal O_X^m\xrightarrow{J}\mathcal O_X^n\longrightarrow\Omega_{X/\Bbbk}\longrightarrow0,\qquad J=\left(\frac{\partial f_j}{\partial x_i}\right),
\]
而 $\operatorname{Sing}(X)$ 由 $J$ 的所有 $(n-d)\times(n-d)$ 子式定义.

\subsection{局部完全交与 syntomic 态射}\label{src:A.3.5-subsection}
\begin{definition}\label{src:A.3.15}
域 $\Bbbk$ 上的局部有限型概形 $X$ 在 $p\in X$ 处为\newterm{局部完全交}{local complete intersection} (或在 $p$ 处为 lci), 如果存在仿射开邻域 $p\in\Spec A\subseteq X$, 使得 $A$ 为 $\Bbbk$ 上的全局完全交, 即 $A\simeq\Bbbk[x_1,\ldots,x_n]/(f_1,\ldots,f_c)$, 且 $\dim A=n-c$. 称概形 $X$ 为局部完全交, 如果它在每个点处都是局部完全交.
\end{definition}
\begin{proposition}\label{src:A.3.16}
对域 $\Bbbk$ 上的局部有限型概形 $X$ 和点 $p\in X$, 下列条件等价:
\begin{enumerate}[(1)]
\item $X$ 在 $p$ 处为局部完全交;
\item 局部环 $\mathcal O_{X,x}\simeq R/(f_1,\ldots,f_c)$, 其中 $R$ 为正则局部环, $f_1,\ldots,f_c\in R$ 为正则序列;
\item 完备化 $\widehat{\mathcal O}_{X,x}\simeq R/(f_1,\ldots,f_c)$, 其中 $R$ 为正则完备局部环, $f_1,\ldots,f_c\in R$ 为正则序列.
\end{enumerate}
\end{proposition}
\begin{proof}
见 \cite[Tags 00S8, 09PY]{SP}.
\end{proof}
对域 $\Bbbk$ 上的局部有限型概形, 有如下蕴含:
\[
\text{光滑}\Longrightarrow\text{局部完全交}\Longrightarrow\text{Gorenstein}\Longrightarrow\text{Cohen--Macaulay}.
\]
相对概念如下:
\begin{definition}\label{src:A.3.17}
称概形态射 $f:X\to S$ 为 \newterm{syntomic 态射}{syntomic} (或\textit{平坦局部完全交态射}), 如果 $f$ 为 fppf 态射且每条纤维都是局部完全交. 称 $f:X\to S$ 在 $x\in X$ 处为 syntomic 态射, 如果存在 $x$ 的开邻域 $U$, 使得 $f|_U$ 为 syntomic 态射.
\end{definition}
syntomic 态射具有类似于局部完全交的局部结构.
\begin{proposition}\label{src:A.3.18}
态射 $f:X\to S$ 在 $x\in X$ 处为 syntomic 态射, 当且仅当存在仿射开邻域 $x\in\Spec A\subseteq X$ 和 $\Spec B\subseteq Y$, 满足 $f(\Spec A)\subseteq\Spec B$, 使得 $A\simeq B[x_1,\ldots,x_n]/(f_1,\ldots,f_c)$, 且 $\Spec A\to\Spec B$ 的每条非空纤维的维数为 $n-c$.
\end{proposition}
\begin{proof}
见 \cite[IV$_4$ \S$19.3$]{EGA}, \cite[VIII \S$1$]{SGA6} 和 \cite[Tag 01UB]{SP}.
\end{proof}
